\documentclass[10pt,a4paper]{amsart}
\usepackage{amsmath,amssymb,amsthm,graphicx}
\usepackage{verbatim}
\newenvironment{reference}{\comment}{\endcomment}
%\newenvironment{reference}{}{}
\newenvironment{slogan}{\comment}{\endcomment}
\newenvironment{history}{\comment}{\endcomment}

% For commutative diagrams we use Xy-pic
\usepackage[all]{xy}

% We use 2cell for 2-commutative diagrams.
\xyoption{2cell}
% Theorem environments.
%
\theoremstyle{plain}
\newtheorem{theorem}[subsection]{Theorem}
\newtheorem{proposition}[subsection]{Proposition}
\newtheorem{lemma}[subsection]{Lemma}

\theoremstyle{definition}
\newtheorem{definition}[subsection]{Definition}
\newtheorem{example}[subsection]{Example}
\newtheorem{exercise}[subsection]{Exercise}
\newtheorem{situation}[subsection]{Situation}

\theoremstyle{remark}
\newtheorem{remark}[subsection]{Remark}
\newtheorem{remarks}[subsection]{Remarks}

\numberwithin{equation}{subsection}

% Macros
%
\def\lim{\mathop{\mathrm{lim}}\nolimits}
\def\colim{\mathop{\mathrm{colim}}\nolimits}
\def\Spec{\mathop{\mathrm{Spec}}}
\def\Hom{\mathop{\mathrm{Hom}}\nolimits}
\def\Ext{\mathop{\mathrm{Ext}}\nolimits}
\def\SheafHom{\mathop{\mathcal{H}\!\mathit{om}}\nolimits}
\def\SheafExt{\mathop{\mathcal{E}\!\mathit{xt}}\nolimits}
\def\Sch{\mathit{Sch}}
\def\Mor{\mathop{\mathrm{Mor}}\nolimits}
\def\Ob{\mathop{\mathrm{Ob}}\nolimits}
\def\Sh{\mathop{\mathit{Sh}}\nolimits}
\def\NL{\mathop{N\!L}\nolimits}
\def\CH{\mathop{\mathrm{CH}}\nolimits}
\def\proetale{{pro\text{-}\acute{e}tale}}
\def\etale{{\acute{e}tale}}
\def\QCoh{\mathit{QCoh}}
\def\Ker{\mathop{\mathrm{Ker}}}
\def\Im{\mathop{\mathrm{Im}}}
\def\Coker{\mathop{\mathrm{Coker}}}
\def\Coim{\mathop{\mathrm{Coim}}}

% Boxtimes
%
\DeclareMathSymbol{\boxtimes}{\mathbin}{AMSa}{"02}

%
% Macros for moduli stacks/spaces
%
\def\QCohstack{\mathcal{QC}\!\mathit{oh}}
\def\Cohstack{\mathcal{C}\!\mathit{oh}}
\def\Spacesstack{\mathcal{S}\!\mathit{paces}}
\def\Quotfunctor{\mathrm{Quot}}
\def\Hilbfunctor{\mathrm{Hilb}}
\def\Curvesstack{\mathcal{C}\!\mathit{urves}}
\def\Polarizedstack{\mathcal{P}\!\mathit{olarized}}
\def\Complexesstack{\mathcal{C}\!\mathit{omplexes}}
% \Pic is the operator that assigns to X its picard group, usage \Pic(X)
% \Picardstack_{X/B} denotes the Picard stack of X over B
% \Picardfunctor_{X/B} denotes the Picard functor of X over B
\def\Pic{\mathop{\mathrm{Pic}}\nolimits}
\def\Picardstack{\mathcal{P}\!\mathit{ic}}
\def\Picardfunctor{\mathrm{Pic}}
\def\Deformationcategory{\mathcal{D}\!\mathit{ef}}
\usepackage{cleveref}
\crefname{originalsection}{original section}{original sections}
\Crefname{originalsection}{Original section}{Original sections}
\makeatletter
\newlabel{section-axioms}{{07XJ}{source}{}{sourcepointer-1}{}}
\newlabel{section-axioms@cref}{{[originalsection][0][]07XJ}{[1][0][]source}{}{}{}}
\makeatother
\title{Behaviour of the degree map under products}
\author{The Stacks Project Authors}
\date{stacks-project:weil.tex:local-snapshot-20261012; GNU FDL 1.2 or later}
\begin{document}
\maketitle
\noindent\textit{Editorial note.} All delivered mathematical text is the author's own, copied from the pinned Stacks Project chapter file weil.tex; no statement or proof step was rewritten, reordered or removed.\par\medskip
\noindent\textit{Editorial note.} The counted proof is the author's complete argument: the derivation of the additive part from the Kunneth decomposition, and the computation of the lowest Kunneth component of the class of the diagonal of a product.\par\medskip
\noindent\textit{Editorial note.} The label lemma-trace-product names the counted statement, and the proof environment immediately following it in the source is delivered verbatim as its proof.\par\medskip
\noindent\textit{Editorial note.} The counted statement is carried by the author's own lemma environment in the chapter source.\par\medskip
\noindent\textit{Editorial note.} Pre-open gate: python operations/agent-authoring/tools/precard\_gate.py --source weil.tex --label lemma-trace-product --start 4946 --end 4991 --local-id 28130 returned RESULT=OK (L28130 free, occupancy FREE, public interval clear). Pre-assembly tag self-check: canonical tag resolved and not present in reference/public-items.jsonl.\par\medskip
\noindent\textit{Editorial note.} The statement and proof are copied from the pinned chapter file; the cited inputs are kept as external pointers with their canonical Stacks tags.\par\medskip
\noindent\textit{Allowed standard prerequisites (editorial).}\begin{itemize}
\item The Kunneth decomposition of the cohomology of a product. Axiom (A5) identifying the cohomology of a product with the tensor product of the cohomologies. Provides the tensor decomposition in which the components of the class of the diagonal are written.
\item The class of the diagonal of a product as the intersection product of the pullbacks of the diagonals. The identification of the diagonal of a product with the intersection product of the pulled back diagonals, together with the compatibility of the cycle class map with intersection products. Expresses the class of the diagonal of the product through the classes of the diagonals of the factors.
\item Definition of the degree map through pairing with the class of the diagonal. Axiom (A6), defining the degree map as the pairing with the class of the diagonal. Shows that the pairing of the lowest Kunneth component of the diagonal with the degree maps of the factors is one.
\end{itemize}\medskip
\ifdefined\hypertarget\hypertarget{sourcepointer-1}{}\fi
\noindent\textit{Source reference.} Original section 07XJ: sources/stacks-weil/pinned-original-source.tex, source lines 2568--2568. Prerequisite pointer invoked at proof line(s) 4949; it names a standard fact of the same chapter that the counted argument uses, and every step of the counted proof built on it is retained verbatim in the delivered body.\par\medskip
\begin{lemma}
\label{lemma-trace-product}
Assume given data (D0), (D1), and (D2') satisfying axioms (A1) -- (A7).
Then axiom (B) of Section \ref{section-axioms} holds.
\end{lemma}

\begin{proof}
Axiom (B)(a) is immediate from axiom (A5).
Let $X$ and $Y$ be nonempty smooth projective schemes over $k$
equidimensional of dimensions $d$ and $e$. To see that axiom (B)(b)
holds, observe that the diagonal $\Delta_{X \times Y}$ of $X \times Y$
is the intersection product of the pullbacks of the diagonals
$\Delta_X$ of $X$ and $\Delta_Y$ of $Y$ by the projections
$p : X \times Y \times X \times Y \to X \times X$ and
$q : X \times Y \times X \times Y \to Y \times Y$.
Compatibility of $\gamma$ with intersection products then gives
that
$$
\gamma([\Delta_{X \times Y}]) \in
H^{2d + 2e}(X \times Y \times X \times Y)(d + e)
$$
is the cup product of the pullbacks of $\gamma([\Delta_X])$
and $\gamma([\Delta_Y])$ by $p$ and $q$. Write
$$
\gamma([\Delta_{X \times Y}]) = \sum \eta_{X \times Y, i}
\text{ with }
\eta_{X \times Y, i} \in
H^i(X \times Y) \otimes H^{2d + 2e - i}(X \times Y)(d + e)
$$
and similarly $\gamma([\Delta_X]) = \sum \eta_{X, i}$ and
$\gamma([\Delta_Y]) = \sum \eta_{Y, i}$. The observation above
implies we have
$$
\eta_{X \times Y, 0} =
\sum\nolimits_{i \in \mathbf{Z}} p^*\eta_{X, i} \cup q^*\eta_{Y, -i}
$$
(If our cohomology theory vanishes in negative degrees, which will
be true in almost all cases, then only the term for $i = 0$ contributes
and $\eta_{X \times Y, 0}$ lies in
$H^0(X) \otimes H^0(Y) \otimes H^{2d}(X)(d) \otimes H^{2e}(Y)(e)$ as expected,
but we don't need this.) Since $\lambda_X : H^{2d}(X)(d) \to F$ and 
$\lambda_Y : H^{2e}(Y)(e) \to F$ send $\eta_{X, 0}$ and $\eta_{Y, 0}$
to $1$ in $H^*(X)$ and $H^*(Y)$, we see that $\lambda_X \otimes \lambda_Y$
sends $\eta_{X \times Y, 0}$ to $1$ in
$H^*(X) \otimes H^*(Y) = H^*(X \times Y)$ and the proof is complete.
\end{proof}

\end{document}
